\section*{TÓM TẮT ĐỒ ÁN}
\thispagestyle{empty}

Đồ án này tập trung nghiên cứu, thiết kế và chế tạo một mô hình thiết bị đeo tay thông minh (Smart Wearable) giám sát sức khỏe thời gian thực tích hợp công nghệ Internet vạn vật (IoT). Hệ thống có khả năng thu thập liên tục các chỉ số sinh hiệu cốt lõi của cơ thể người bao gồm nhịp tim, nồng độ bão hòa oxy trong máu ngoại vi ($SpO_2$), nhiệt độ da và độ ẩm cơ thể, sau đó đồng bộ trực tiếp lên máy chủ trung tâm để phục vụ mục đích giám sát y tế từ xa (IoMT).

Về mặt phần cứng, bo mạch đeo tay được thiết kế nhỏ gọn dạng bát giác 2 lớp trên phần mềm Altium Designer với đường kính 42mm phù hợp khung vỏ đồng hồ đeo tay. Hệ thống sử dụng vi điều khiển thế hệ mới \textbf{ESP32-C3} kiến trúc RISC-V hiệu năng cao và tiết kiệm điện, tích hợp cảm biến quang học sinh trắc học \textbf{MAX30102} qua giao thức bus I2C để đo nhịp tim/$SpO_2$, cảm biến kỹ thuật số siêu nhỏ \textbf{HDC2022} qua bus I2C để đo nhiệt độ da và độ ẩm cơ thể, hiển thị cục bộ tại chỗ trên màn hình \textbf{OLED SSD1306 1.3 inch}. Khối nguồn của thiết bị sử dụng pin sạc Lithium 3.7V tích hợp mạch quản lý sạc an toàn tuyến tính \textbf{TP4056} phối hợp IC bảo vệ chuyên dụng \textbf{DW01A} và Mosfet FS8205A để ngăn ngừa các sự cố quá dòng, quá xả hay chập mạch.

Về phần mềm nhúng (Firmware), vi điều khiển được lập trình bộ xử lý tín hiệu số (DSP) bao gồm bộ lọc khử trôi dòng DC (EMA Filter), bộ lọc thông thấp số IIR để triệt tiêu nhiễu cơ học, và giải thuật \textbf{Phát hiện Đỉnh thích ứng động (Adaptive Peak Detection)} để xác định chính xác nhịp mạch đập. Đặc biệt, để giải quyết vấn đề sụt giảm biên độ tín hiệu quang học PPG khi đo tại cổ tay, em đã phát triển thuật toán \textbf{Bù sai số tĩnh thích ứng đo ở cổ tay (Wrist Offset Compensation)}, giúp tự động co dãn tuyến tính dữ liệu tính toán $SpO_2$ về đúng dải sinh lý chuẩn. Thiết bị cũng tích hợp cơ chế cấu hình Wi-Fi và máy chủ API không dây thông minh thông qua \textbf{AP Mode \& Captive Portal} tại địa chỉ \texttt{192.168.4.1}, lưu cấu hình vào bộ nhớ Flash qua thư viện \texttt{Preferences} và giới hạn công suất phát sóng RF Wi-Fi ở mức \texttt{11dBm} để loại bỏ lỗi sụt dòng gây reset chip.

Về phía máy chủ (Server), em đã phát triển hệ thống API trung tâm có khả năng chạy song song trên cả nền tảng \textbf{Node.js} và \textbf{PHP}, thực hiện lưu trữ dữ liệu dạng tệp phẳng JSON (\texttt{iot\_store.json}) với giải pháp ghi tệp không đồng bộ và giới hạn lịch sử 120 điểm đo gần nhất. Giao diện Dashboard Web y tế chuyên dụng được xây dựng theo phong cách Dark Mode hiện đại hiển thị danh sách thiết bị hoạt động, tự động cập nhật dữ liệu sinh hiệu thời gian thực thông qua Ajax Polling định kỳ mỗi 3 giây, trực quan hóa biểu đồ lịch sử đo bằng thư viện Chart.js và tích hợp cơ chế phân quyền quản trị Bác sĩ/Y tá thông qua Session Cookie. Thử nghiệm thực tế cho thấy hệ thống hoạt động rất ổn định, đạt sai số tuyệt đối trung bình (MAE) thấp chỉ $\approx 1.7\text{ BPM}$ đối với nhịp tim và $\approx 0.34\%$ đối với $SpO_2$ so với thiết bị y tế Beurer PO30 chuyên dụng.

\vspace{6pt}
\noindent\textbf{Từ khóa:} Smart Wearable, ESP32-C3, MAX30102, HDC2022, PPG, SpO2, IoT, API Server, Node.js, Web Dashboard, lọc số DSP.
\cleardoublepage
